\medskip

L'espace est muni d'un repère orthonormée \Oijk.

On considère:
\begin{itemize}
\item[$\bullet~~$]$d_1$ la droite passant par le point H(2~;~3~;~0) et de vecteur directeur $\vect{u}\begin{pmatrix}1\\-1\\1\end{pmatrix}$ ;
\item[$\bullet~~$]$d_2$  la droite de représentation paramétrique :
\[\left\{\begin{array}{l c l}
x &=& 2k - 3\\
y &=&k\\
z &=&5
\end{array}\right. \quad \text{où }\:k \:\text{décrit }\: \R.\]
\end{itemize}

\medskip

Le but de cet exercice est de déterminer une représentation paramétrique d'une droite $\Delta$ qui soit perpendiculaire aux droites $d_1$ et $d_2$.

\medskip

\begin{enumerate}
\item
	\begin{enumerate}
		\item Déterminer un vecteur directeur $\vect{v}$ de la droite $d_2$. 
		\item Démontrer que les droites $d_1$ et $d_2$ ne sont pas parallèles. 
		\item Démontrer que les droites $d_1$ et $d_2$  ne sont pas sécantes. 
		\item Quelle est la position relative des droites $d_1$ et $d_2$ ?
	\end{enumerate}
\item 
	\begin{enumerate}
		\item Vérifier que le vecteur $\vect{w}\begin{pmatrix}-1\\2\\3\end{pmatrix}$ est orthogonal à $\vect{u}$ et à $\vect{v}$.
		\item On considère le plan $P$ passant par le point H et dirigé par les vecteurs $\vect{u}$ et $\vect{w}$.
		
On admet qu'une équation cartésienne de ce plan est :
\[5x + 4y - z - 22 = 0.\]

Démontrer que l'intersection du plan $P$ et de la droite $d_2$ est le point M(3~;~3~;~5).
	\end{enumerate}
\item Soit $\Delta$ la droite de vecteur directeur $\vect{w}$ passant par le point M. 

Une représentation paramétrique de $\Delta$ est donc donnée par:

\[\left\{\begin{array}{l c l}
x &=&- r + 3 \\
y &= &2r + 3\\
z &=& 3r + 5
\end{array}\right. \quad \text{où }\:r\: \text{décrit }\R.\]

	\begin{enumerate}
		\item Justifier que les droites $\Delta$ et $d_1$ sont perpendiculaires
en un point L dont on déterminera les coordonnées.
		\item Expliquer pourquoi la droite $\Delta$ est solution du problème posé.
	\end{enumerate}
\end{enumerate}

\bigskip
